\documentclass[11pt]{article}
\usepackage[margin=2.4cm]{geometry}
\usepackage{amsmath,amssymb,booktabs}
\usepackage[hidelinks]{hyperref}

\newcommand{\dd}{\mathrm{d}}
\newcommand{\T}{^{\mathsf T}}

\title{Stage-consistent motor command in the RK4 step\\of the flexible-joint arm}
\author{}
\date{}

\begin{document}
\maketitle
\vspace{-3em}

\section{What was wrong, and what the change fixes}

\paragraph{Symptom.} In air --- no cutting force, gravity off, the arm initialised exactly on its
tracking equilibrium --- the TCP settled at a steady offset behind the command, along the feed.
Its size scaled linearly with the step $h$, so it was numerical, not physical:
\[
  \delta p_x\;\approx\;-0.63\,v\,h
  \qquad(h=109.4\,\mu\text{s}):\qquad
  \begin{array}{c|cccc}
    v\ [\text{mm/s}] & 3 & 40 & 120 & 240\\\hline
    \delta p_x\ [\mu\text{m}] & -0.2 & -2.7 & -8.2 & -16.4
  \end{array}
\]

\paragraph{Cause.} The integrator held the motor command constant over each RK4 step. That is
equivalent to driving the joint springs with $\theta_f(t-\tfrac h2)$: a pure half-step delay
(\S\ref{sec:hold}). A delay of $55\,\mu$s is negligible as a \emph{phase} at the structural modes
($0.45^\circ$ at 23\,Hz), but as a \emph{position} it is $v\,h/2$ --- and the command moves fast
relative to the deflections being measured.

\paragraph{Why it mattered here.}
\begin{itemize}
  \item The quantity under study, the cut-induced deflection, is $1$--$20\,\mu$m: the same order as
        the artefact.
  \item At fixed chip load the feed spans $80\times$ across the tooth-passing sweep
        ($3$--$240$\,mm/s), so the artefact was not a constant offset but a \emph{trend along the
        sweep axis} ($0.2$ to $16\,\mu$m), indistinguishable from a physical effect of rpm and
        tooth count.
  \item The scheme was first order in $h$ regardless of RK4, so refining $h$ moved the answer.
\end{itemize}

\paragraph{Effect on the model/simulation comparison.} 10000\,rpm, 4 teeth, $v=120$\,mm/s,
engaged-span DC deviation [$\mu$m]:
\begin{center}
\begin{tabular}{lrrr}
\toprule
 & hold & stage-consistent & linear model \\
\midrule
$x$ & 2.96  & 11.27 & 11.34 \\
$y$ & $-0.16$ & $-0.14$ & $-0.24$ \\
$z$ & 16.97 & 18.84 & 17.04 \\
\bottomrule
\end{tabular}
\end{center}
\begin{itemize}
  \item $x$: the apparent $8.4\,\mu$m model error was the integrator. After the change the linear
        model and the simulation agree to $0.07\,\mu$m.
  \item $z$: the hold also shifted $z$ (by $-1.87\,\mu$m, \S\ref{sec:offset}) and happened to land
        it on the model, hiding a real $1.8\,\mu$m ($10\%$) discrepancy that is now visible.
  \item Forces are unchanged ($-97.13$\,N in both), so every force result stands.
\end{itemize}

\paragraph{Result.} Evaluating the command at each stage's own time removes the delay. The
in-air offset becomes $-0.11\,\mu$m \emph{at every} $h$ --- converged, and physical
(\S\ref{sec:new}). The input error drops from $\mathcal O(h)$ to $\mathcal O(h^2)$ at no extra
cost. The rest of this note derives both statements.

\section{Model}

Joints split into flexible and rigid, $q=(q_f,q_r)\in\mathbb R^{n_f}\times\mathbb R^{n_r}$.
The motor command $\theta(t)\in\mathbb R^{n}$ is a known function of time, available on the
simulation grid only as samples $\theta_n=\theta(t_n)$, $t_n=nh$. Rigid joints are prescribed,
$q_r\equiv\theta_r$. The flexible joints obey
\begin{equation}
  \ddot q_f=\Phi\big(q,\dot q,\tau_u+\tau_{\mathrm{ext}}\big),\qquad
  \tau_u=-K\,(q_f-\theta_f)-D\,(\dot q_f-\dot\theta_f),
  \label{eq:model}
\end{equation}
with $\Phi$ the forward dynamics under the prescribed rigid joints, $K,D$ diagonal, and
$\tau_{\mathrm{ext}}=J_f\T f_{\mathrm{ext}}$. With $x=(q_f,\dot q_f)$ and $u=(\theta,\dot\theta)$,
\begin{equation}
  \dot x=f\big(x,u(t),\tau_{\mathrm{ext}}(t)\big)=:F(t,x).
\end{equation}
The system is non-autonomous only through the command and the external torque.

\section{Classical RK4}

For $\dot x=F(t,x)$,
\begin{align*}
  k_1&=F(t_n,\;x_n), &
  k_2&=F(t_n+\tfrac h2,\;x_n+\tfrac h2k_1),\\
  k_3&=F(t_n+\tfrac h2,\;x_n+\tfrac h2k_2), &
  k_4&=F(t_n+h,\;x_n+hk_3),
\end{align*}
\begin{equation}
  x_{n+1}=x_n+\tfrac h6\,(k_1+2k_2+2k_3+k_4).
\end{equation}
Nodes $c=(0,\tfrac12,\tfrac12,1)$. Local error $\mathcal O(h^5)$, global $\mathcal O(h^4)$ for
$F\in C^4$ --- \emph{provided} stage $i$ evaluates $F$ at $t_n+c_ih$, including its explicit
time dependence.

\section{Previous scheme: zero-order hold}

Every stage used the command at the start of the step:
\begin{equation}
  k_i=f\big(x^{(i)},\,u_n,\,\tau_n\big),\qquad i=1,\dots,4.
\end{equation}
This is classical RK4 applied to the \emph{held} system
\begin{equation}
  \dot x=f\big(x,u_h(t),\tau_h(t)\big),\qquad u_h(t)=u_n,\;\; t\in[t_n,t_{n+1}).
\end{equation}
RK4 solves the held system to $\mathcal O(h^4)$; the error is therefore dominated by the model
perturbation $u_h-u$.

\subsection{The hold is a half-step delay}\label{sec:hold}

For smooth $u$ and $t\in[t_n,t_{n+1})$,
\begin{equation}
  u_h(t)-u(t)=-\dot u(t_n)\,(t-t_n)+\mathcal O(h^2),
\end{equation}
whose step mean is $-\tfrac h2\dot u+\mathcal O(h^2)$; the remainder is a zero-mean sawtooth at
$1/h$. In the frequency domain the hold is
\begin{equation}
  H(i\omega)=\frac{1-e^{-i\omega h}}{i\omega h}
  =e^{-i\omega h/2}\,\operatorname{sinc}\!\big(\tfrac{\omega h}{2}\big).
\end{equation}
In the structural band $\omega\le2\pi\cdot23$\,Hz with $h\approx10^{-4}$\,s, $\omega h\approx0.015$:
unit gain to $\mathcal O((\omega h)^2)$ and a pure delay $h/2$. The flexible joints are driven by
$\theta_f(t-\tfrac h2)$, and the scheme is \textbf{first order}, whatever the order of the
integrator.

\subsection{Tracking offset}\label{sec:offset}

On a constant-velocity path the step-averaged spring force vanishes, so the hold contributes
$q_f(t)=\theta_f(t-\tfrac h2)$ on top of the $h$-independent physical deflection. The rigid joints
are not integrated and carry no delay. The TCP error is therefore
\begin{equation}
  \delta p\;\approx\;-\tfrac h2\,J_f(\theta)\,\dot\theta_f,
  \qquad\text{not}\qquad -\tfrac h2\,J\dot\theta=-\tfrac h2\,v .
  \label{eq:offset}
\end{equation}
On this arm the rigid wrist moves against the feed, so $J_f\dot\theta_f$ exceeds $v$.
At 3333\,rpm, 4 teeth, $v=40$\,mm/s along workpiece $x$ (values in mm/s):
\begin{center}
\begin{tabular}{lrrr}
\toprule
 & $x$ & $y$ & $z$ \\
\midrule
$J\dot\theta$            & 40.00 & 0.00 & 0.00 \\
$J_f\dot\theta_f$        & 50.06 & 0.00 & 10.78 \\
$J_r\dot\theta_r$        & $-10.06$ & 0.00 & $-10.78$ \\
\bottomrule
\end{tabular}
\end{center}
Eq.~\eqref{eq:offset} predicts $\delta p_x/(vh)=\tfrac12\cdot50.06/40.00=0.626$ and
$\delta p_z/\delta p_x=10.78/50.06=0.215$.

\subsection{Measured}

Pure tracking in air --- no cutting force, gravity off, $v=40$\,mm/s --- steady feed-direction
error:
\begin{center}
\begin{tabular}{rrrr}
\toprule
$h$ [$\mu$s] & $\delta p_x$ hold [$\mu$m] & $\delta p_x$ staged [$\mu$m] &
$(\delta p_x^{\mathrm{hold}}-\delta p_x^{\mathrm{staged}})/(vh)$ \\
\midrule
109.38 & $-2.859$ & $-0.1105$ & 0.628 \\
 54.69 & $-1.485$ & $-0.1105$ & 0.628 \\
 27.34 & $-0.798$ & $-0.1105$ & 0.629 \\
\bottomrule
\end{tabular}
\end{center}
Exact first-order scaling, coefficient $0.628$ against $0.626$ predicted. In a full coupled pass
at 120\,mm/s, switching scheme shifted the engaged-span DC deviation by $+8.31\,\mu$m in $x$ and
$+1.87\,\mu$m in $z$ (predicted $8.21$ and $1.77$; ratio $0.225$ against $0.215$).

\section{New scheme: stage-consistent command}\label{sec:new}

Given the samples at both ends of the step, each stage takes the command at its own node:
\begin{align}
  k_1&=f\big(x_n,\,u_n,\,\tau_n\big), &
  k_2&=f\big(x_n+\tfrac h2k_1,\,\bar u_n,\,\tau_n\big),\nonumber\\
  k_3&=f\big(x_n+\tfrac h2k_2,\,\bar u_n,\,\tau_n\big), &
  k_4&=f\big(x_n+hk_3,\,u_{n+1},\,\tau_n\big),
\end{align}
with $\bar u_n=\tfrac12(u_n+u_{n+1})$, applied to both $\theta$ and $\dot\theta$.

\paragraph{Equivalence.} Let $u_\ell$ be the piecewise-linear interpolant of the samples. Then
$u_\ell(t_n+c_ih)=(u_n,\bar u_n,\bar u_n,u_{n+1})$, so the scheme is classical RK4 applied to
\begin{equation}
  \dot x=f\big(x,u_\ell(t),\tau_h(t)\big).
\end{equation}

\paragraph{Error.} On $[t_n,t_{n+1}]$,
\begin{equation}
  u_\ell(t)-u(t)=-\tfrac12\,\ddot u(\xi)\,(t-t_n)(t_{n+1}-t)=\mathcal O(h^2),
\end{equation}
step mean $-\tfrac{h^2}{12}\ddot u$. There is no $\mathcal O(h)$ term and no delay: the
interpolant is exact for affine $u$. The input error is $\mathcal O(h^2)$ globally, the
integration $\mathcal O(h^4)$; overall $\mathcal O(h^2)$.

On a constant-speed straight path $\ddot\theta\neq0$ only because the Jacobian rotates, so the
residual is below measurement: $\delta p_x=-0.1105\,\mu$m at every $h$ in the table.

\paragraph{The remaining offset is physical.} At steady tracking $\dot q\approx\dot\theta$,
$\ddot q\approx\ddot\theta$, and~\eqref{eq:model} with gravity off gives, for
$e=q_f-\theta_f$,
\begin{equation}
  K\,e\;\approx\;-\big[M(q)\,\ddot\theta+b(q,\dot\theta)\big]_f,
\end{equation}
independent of $h$: $0.11\,\mu$m here.

\section{What is still held}

$\tau_{\mathrm{ext}}$ is sampled once per step by the process engine and held, so the cutting
force enters the loop with the same delay $h/2$ as before. Phase lag $\omega h/2$ at
$h=109.4\,\mu$s:
\begin{center}
\begin{tabular}{lrr}
\toprule
 & $f$ [Hz] & $\omega h/2$ \\
\midrule
highest structural mode & 23   & $0.45^\circ$ \\
tooth passing, 3333\,rpm $\times$ 4 & 222 & $4.4^\circ$ \\
tooth passing, 10000\,rpm $\times$ 8 & 1333 & $26^\circ$ \\
\bottomrule
\end{tabular}
\end{center}
Negligible where the regenerative loop lives; a phase shift of the forcing itself at the fastest
cells, where the forcing lies far above the structural band.

\section{Recording convention}

Sample $n$ stores $(t_n,\theta_n,q_n,f_{\mathrm{ext},n})$ taken \emph{before} step $n$, so
$\delta p_n=p(q_n)-p_{\mathrm{cmd}}(t_n)$ compares quantities at one instant. The former recorder
stored the post-step state against $t_n$; since the post-step assembly advances only the
integrated (flexible) joints, that added $+h\,J_f\dot\theta_f$ to every $\delta p$ --- the same
direction as~\eqref{eq:offset}, opposite sign, twice the size. It was removed separately and is
not part of the offsets above.

\section{Implementation}

\begin{center}
\begin{tabular}{lp{0.55\textwidth}}
\toprule
\texttt{robotsim/solver.py}   & \texttt{\_rk4\_step}: stage commands $(u_n,\bar u_n,\bar u_n,u_{n+1})$
                                 when \texttt{theta\_end} is given \\
\texttt{robotsim/dynamics.py} & \texttt{Simulator.step}: forwards \texttt{theta\_end},
                                 \texttt{thetaD\_end}; records at $t_n$ \\
\texttt{analysis/sim\_coupled.py} & loop passes $\theta_{n+1}$, $\dot\theta_{n+1}$
                                 (clamped at the last sample) \\
\bottomrule
\end{tabular}
\end{center}
Without \texttt{theta\_end} the step reduces to the zero-order hold, so callers that do not pass
it are unchanged.

\end{document}
